\documentclass[12pt]{article}
\usepackage[utf8]{inputenc}
\usepackage[T1]{fontenc}
\usepackage[margin=1in]{geometry}
\usepackage{setspace}
\usepackage{xcolor}
\usepackage[colorlinks=true, urlcolor=blue!60!black, linkcolor=blue!60!black]{hyperref}
\setlength{\parindent}{0em}
\setlength{\parskip}{0.6em}
\renewcommand{\baselinestretch}{1.0}
\begin{document}
\begin{center}
{\LARGE \textit{National Pork Producers Council v. Ross}, 598 U.S. 356 (2023)\par}
\vspace{0.5em}
{\large Case Brief --- Two-Page Edition\par}
\end{center}

\section{Memory Jogger}
The Court unanimously rejected an ``almost per se'' extraterritoriality rule and, 5--4, upheld California's Proposition 12 ban on in-state sales of pork from cruelly confined pigs (\href{https://www.law.cornell.edu/supremecourt/text/21-468}{598 U.S.\ 356}).

\section{Detailed Case Facts}
Proposition 12 barred in-state sales of whole pork from breeding pigs ``confined in a cruel manner.'' Pork-industry groups alleged that California buys about 13\% of the nation's pork but imports nearly all of it, so compliance costs would force producers nationwide to restructure. They conceded the law does not discriminate against out-of-state producers.

\section{Procedural History}
The district court dismissed under Rule 12(b)(6); the Ninth Circuit affirmed.

\section{Judicial Votes}
Affirmed, 5--4. Gorsuch wrote for Thomas, Sotomayor, Kagan, and Barrett in key parts, with shifting pluralities on \textit{Pike}; Sotomayor and Barrett concurred in part; Roberts (joined by Alito, Kavanaugh, Jackson) and Kavanaugh dissented in part.

\section{Holding}
Antidiscrimination lies at the ``very core'' of the Dormant Commerce Clause, and there is no ``almost per se'' rule against nondiscriminatory laws with out-of-state effects. \textit{Pike} survives but mainly smokes out protectionism or protects interstate transportation; courts must use ``extreme caution'' absent congressional action.

\section{Analysis of Opinions}
Gorsuch read \textit{Baldwin}, \textit{Brown-Forman}, and \textit{Healy} as protectionist price-control cases and warned that petitioners' rule would imperil ordinary state laws. A three-Justice plurality found costs and moral benefits incommensurable; a four-Justice plurality found no substantial burden pleaded. Sotomayor and Kagan preserved \textit{Pike} but found no burden; Barrett found a burden but no judicial scale. Roberts would have remanded for full balancing of market-wide harms, and Kavanaugh warned of ``California knows best'' balkanization.

\section{Examples: Future Applications}
\textbf{Same-Side:} a neutral ban on selling fur or foie gras in the State is upheld. \textbf{Opposite-Side:} a law tying in-state drug prices to out-of-state prices falls under \textit{Healy}. \textbf{Fence-Sitter:} sales bans tied to out-of-state firms' labor or reproductive-health policies.

\section{Critique}
Originalists praised the refusal to expand an atextual doctrine; others fear balkanization by large States. Progressives hailed state autonomy, while doctrinal commentators note that the fractured vote leaves \textit{Pike} unsettled.

\section{Key Quotations}
\begin{quote}
``Companies that choose to sell products in various States must normally comply with the laws of those various States. Assuredly, under this Court's dormant Commerce Clause decisions, no State may use its laws to discriminate purposefully against out-of-state economic interests. But the pork producers do not suggest that California's law offends this principle. Instead, they invite us to fashion two new and more aggressive constitutional restrictions on the ability of States to regulate goods sold within their borders. We decline that invitation. While the Constitution addresses many weighty issues, the type of pork chops California merchants may sell is not on that list.''
\end{quote}
\noindent --- Justice Gorsuch, Opinion of the Court, 598 U.S.\ at 364.

\end{document}
